\input{build/foundations-preamble.tex}
\usepackage[authoryear]{natbib}
\renewcommand{\refname}{સંદર્ભગ્રંથો}
\renewcommand{\partname}{ભાગ}
\newenvironment{intro}{}{}
\newcommand{\Id}[1]{\mathord{\mathrm{Id}_{#1}}}
\newcommand{\equivrep}[2]{[#1]_{#2}}
\newcommand{\equivclass}[2]{#1/_{\!{#2}}}
\newcommand{\Intequiv}{\mathrel{\sim_{\Int}}}
\newcommand{\Ratequiv}{\mathrel{\sim_{\Rat}}}
\newcommand{\Realequiv}{\mathrel{\sim_{\Real}}}
\newcommand{\funrestrictionto}[2]{#1\mathord{\restriction}_{#2}}
\newcommand{\funimage}[2]{#1[#2]}
\newcommand{\emptyseq}{\Lambda}
\newcommand{\liff}{\mathbin{\leftrightarrow}}
\newcommand{\dom}[1]{\mathrm{dom}(#1)}
\newcommand{\ran}[1]{\mathrm{ran}(#1)}
\newcommand{\comp}[2]{#2\circ #1}
\newcommand{\pto}{\mathrel{\ooalign{\hfil$\mapstochar\mkern5mu$\hfil\cr$\to$}}}
\newcommand{\fdefined}{\downarrow}
\newcommand{\fundefined}{\uparrow}
\newcommand{\True}{\mathbb{T}}
\newcommand{\False}{\mathbb{F}}
\newcommand{\lfalse}{\bot}
\newcommand{\ltrue}{\top}
\newcommand{\Obj}[1]{\mathsf{#1}}
\newcommand{\Lang}[1]{\mathcal{#1}}
\NewDocumentCommand{\Frm}{o}{%
  \IfNoValueTF{#1}{\mathrm{Frm}}{\mathrm{Frm}(\Lang{#1})}}
\newcommand{\PVar}{\mathrm{At}_0}
\newcommand{\Struct}[1]{\mathfrak{#1}}
\newcommand{\pAssign}[1]{\mathfrak{#1}}
\newcommand{\pValue}[1]{\overline{\mathfrak{#1}}}
\NewDocumentCommand{\pSat}{t{/} m m}{%
  \pAssign{#2}\IfBooleanTF{#1}{\nvDash}{\vDash}#3}
\newcommand{\Entails}{\vDash}
\newcommand{\ident}{\equiv}
\newcommand{\subst}[2]{#1/#2}
\newcommand{\SSubst}[2]{#1[#2]}
\newcommand{\Subst}[3]{#1[\subst{#2}{#3}]}
\newcommand{\indfrm}{}
\newcommand{\indcase}[3]{\def\indfrm{#1}$#1\ident #2$: #3}
\expandafter\def\csname gutoken@formula\endcsname{સૂત્ર}
\expandafter\def\csname gutoken@valuation\endcsname{સત્યમૂલ્ય-નિયુક્તિ}
\newcommand{\sourcecorrection}[2]{%
  \leavevmode\par\smallskip\noindent
  \fbox{\begin{minipage}{0.94\linewidth}\small
    \textbf{સ્રોત-સુધારો #1.} #2
  \end{minipage}}%
  \par\smallskip}
\definecolor{oldiagcolorD}{HTML}{1973ba}
\newlength{\olphotowidth}
\setlength{\olphotowidth}{.4\textwidth}
\RenewDocumentCommand{\olasset}{O{\olphotowidth} m}{\centering\resizebox{#1}{!}{\input{upstream/#2}}}
\input{build/propositional-available-labels.tex}
\renewcommand{\oliflabeldef}[3]{\ifcsname guavailable@#1\endcsname#2\else#3\fi}
\hypersetup{pdftitle={ગણો અને વિધાનાત્મક તર્કશાસ્ત્ર — ઓપન લોજિક ગુજરાતી}}
\begin{document}
\begin{titlepage}
{\Huge\bfseries ગણો, સંબંધો, વિધેયો, ગણોનું કદ, અંકગણિતીકરણ,
અનંત ગણો અને વિધાનાત્મક તર્કશાસ્ત્ર\par}
\vspace{1em}
{\Large ઓપન લોજિક — ગુજરાતી\par}
\vspace{2em}
ઓપન લોજિક પ્રોજેક્ટના અંગ્રેજી મૂળનો ગુજરાતી અનુવાદ.\par
\vspace{1em}
આ આવૃત્તિમાં ગણો અને વિધેયોના અગાઉના પ્રકરણો સાથે વિધાનાત્મક
તર્કશાસ્ત્રના વાક્યરચના અને અર્થવિચાર પ્રકરણના છ ખંડ ઉમેર્યા છે.
તેમાં સૂત્રોની રચના, એકમાત્ર વાચનીયતા, પ્રતિસ્થાપન, રચના-શ્રેણીઓ,
સત્યમૂલ્ય-નિયુક્તિઓ, સત્યાર્થતા સારણીઓ, સંતોષ, પુનરુક્તિ અને
અર્થાનુસારી ફલિતતા સામેલ છે.\par
\vfill
આ યંત્ર દ્વારા કરેલો અનુવાદ છે. સંપૂર્ણ ગ્રંથનું કામ ચાલુ છે.
આ આવૃત્તિમાં મૂળનાં ૫૯ એકમોનો સમાવેશ છે; કુલ ૭૨૨ એકમો છે.
કેટલીક વિશિષ્ટ પરિભાષા કામચલાઉ છે. ચકાસણીની વિગતો અને
મૂળ પાઠ વિશેની નોંધો સંપાદકીય માહિતીમાં આપેલી છે.\par
\vspace{1em}
\href{https://github.com/OpenLogicProject/OpenLogic}{Open Logic Project}
\quad \href{https://creativecommons.org/licenses/by/4.0/}{CC BY 4.0}\par
\href{https://github.com/KokunoYumeto/OpenLogic-translations}{અનુવાદોનું કેન્દ્ર}\par
\end{titlepage}
\tableofcontents
\clearpage
\input{build/propositional-body.tex}
\clearpage
\section*{સંપાદકીય નોંધો}
\addcontentsline{toc}{section}{સંપાદકીય નોંધો}
આ નોંધો મૂળ પાઠથી અલગ છે. મૂળની ગણિતીય નિશાનીઓ જાળવી છે.
OLFUN-001થી OLFUN-005, OLSIZ-001થી OLSIZ-012, OLARI-001થી
OLARI-004, OLINF-001થી OLINF-002 અને OLPL-001થી OLPL-002 સુધીની
ઓળખવાળી નોંધો સ્થિર અંગ્રેજી મૂળમાં મળેલી ખામીઓ અને ગુજરાતી પાઠમાં
કરેલા મર્યાદિત સુધારા બાજુબાજુ દર્શાવે છે. દરેક સુધારાની ચોક્કસ
જગ્યા, કારણ અને સૂત્રમાં થયેલો ફેરફાર નિર્ણય-અભિલેખમાં નોંધાયેલ છે.

વિધાનાત્મક તર્કશાસ્ત્રના આ વાંચનરૂપમાં સ્થિર મૂળની સામાન્ય ગોઠવણી
અનુસરીને અસત્ય અને સત્યના અચળો તથા પાંચેય તાર્કિક સંયોજકોને મૂળભૂત
ગણ્યા છે. વૈકલ્પિક રીતે વ્યાખ્યાયિત સંયોજકોનો અનુવાદ ગોઠવણી-ટૅગો સાથે
સ્રોત ફાઇલમાં જાળવેલો છે. તે ગોઠવણીમાં પ્રેરણની વ્યાખ્યામાં આવેલા
વધારાના અંતકૌંસને OLPL-001 વડે સુધાર્યો છે. રચના-શ્રેણીની સાબિતીમાં
તાર્કિક સમતુલ્યતાના સંકેતની જગ્યાએ જરૂરી વાક્યરચનાત્મક અભિન્નતાનો
સંકેત OLPL-002 વડે પુનઃસ્થાપિત કર્યો છે.

વિધાનચલ, આણ્વિક સૂત્ર, સત્યમૂલ્ય-નિયુક્તિ, સંતોષ્ય, પુનરુક્તિ અને
નિવાર્ય જેવા પદો ગુજરાતી પાઠ્યપુસ્તક અને વિશ્વકોશના તપાસેલા
તર્કશાસ્ત્રીય પ્રયોગો સાથે સરખાવ્યાં છે. અર્થાનુસારી ફલિતતા અને
સ્થાનિક નિર્ધારણ જેવા સંયોજનો નિષ્ણાત સુધારા માટે ખુલ્લા છે.

અનંત ગણોના પ્રકરણમાં સંવરણની વ્યાખ્યાને સ્વવિધેય ($f\colon A\to A$),
ઉપગણ ($X\subseteq A$) અને આરંભઘટક ($o\in A$) સાથે સ્પષ્ટ કરી છે;
મૂળની સાબિતી અને પછીના બધા ઉપયોગો આ જ પ્રકારધારણા પર આધાર રાખે છે.
સંવરણના સહાયક પરિણામમાં મૂળની ગેરરીતે ગૂંથાયેલી સભ્યસંખ્યા-સમાનતાને
તેની સાબિતી સ્થાપિત કરતી બે સમાનતાઓથી બદલી છે.

ગણનીય માટેનું ગુજરાતી પ્રમાણ ગણનીયને ગણનીય અનંત ગણ માટે વાપરે છે,
જ્યારે Open Logicની આ આવૃત્તિની વ્યાખ્યા સાન્ત અને ખાલી ગણોને પણ
સમાવે છે. બંને પરિગણના-રૂપો પોતાની સ્થાનિક વ્યાખ્યા સાથે જ વાંચવાનાં છે.

પરિભાષા, મૂળ સ્રોતો, પ્રત્યેક અનુવાદિત ખંડની સ્રોત-સંગતિ અને
ચકાસણીની વધુ વિગતો આવૃત્તિની સાથેની ફાઇલોમાં છે.
\begin{thebibliography}{9}
\bibitem[Benacerraf(1965)]{Benacerraf1965}
Benacerraf, Paul. 1965. \emph{What numbers could not be}.
The Philosophical Review 74(1), 47--73.
\bibitem[Cantor(1892)]{Cantor1892}
Cantor, Georg. 1892. Über eine elementare Frage der Mannigfaltigkeitslehre.
\bibitem[Frege(1884)]{Frege1884}
Frege, Gottlob. 1884. \emph{Die Grundlagen der Arithmetik}.
\bibitem[Potter(2004)]{Potter2004}
Potter, Michael. 2004. \emph{Set Theory and Its Philosophy}.
\bibitem[Conway(2006)]{Conway2006}
Conway, John. 2006. \emph{The Power of Mathematics}.
\bibitem[Katz and Katz(2012)]{KatzKatz2012}
Katz, Karin Usadi and Mikhail G. Katz. 2012. Stevin Numbers and Reality.
\bibitem[O'Connor and Robertson(2005)]{OConnorRobertson:RN}
O'Connor, John J. and Edmund F. Robertson. 2005. The real numbers: Stevin to Hilbert.
\bibitem[Hilbert(2013)]{EwaldSieg2013}
Hilbert, David. 2013. \emph{David Hilbert's Lectures on the Foundations of Arithmetic and Logic 1917--1933}.
Edited by William Bragg Ewald and Wilfried Sieg. Springer.
\bibitem[Dedekind(1888)]{Dedekind1888}
Dedekind, Richard. 1888. \emph{Was sind und was sollen die Zahlen?} Vieweg.
\end{thebibliography}
\end{document}
